% docs/discusion_datos.tex
\section{Discusión de datos}

Esta sección analiza los datos y resultados generados con DIAPCE en el contexto del Meta (Colombia), con énfasis en su uso formativo por parte de \textbf{los estudiantes del semillero de investigación y de los laboratorios de materiales de la UCC}. Es relevante precisar que DIAPCE es una herramienta de alcance institucional y académico, no una solución de uso general; por tanto, la discusión se enmarca siempre en el contexto de prácticas, ensayos y ejercicios realizados por los estudiantes dentro del programa formativo de la UCC. La discusión se centra en: (i) la calidad y cobertura del conjunto de datos regional incorporado en el repositorio (\texttt{lib/data/csv\_datos\_1.1.csv}); (ii) la validez de las recomendaciones del sistema para dosificación base y clasificación de aditivos (P0–PP3) con proyecciones a 7, 14 y 28 días; y (iii) la comparación frente a prácticas tradicionales de diseño según ACI 211.1.

\subsection{Conjunto de datos y preprocesamiento}
El conjunto de datos registra parámetros ambientales y de diseño (p. ej., resistencia objetivo, temperatura, humedad, condición de exposición) y variables de mezcla y desempeño. Para su uso en DIAPCE se aplican pasos mínimos de preprocesamiento: depuración de registros incompletos, homogeneización de unidades, detección de valores atípicos y codificación de la etiqueta de aditivo (P0–PP3). Se recomienda documentar las decisiones de limpieza empleadas en cada cohorte de ensayos para garantizar reproducibilidad.

\begin{table}[H]
  \centering
  \caption{Cobertura de rangos del conjunto de datos (1\,872 registros)}
  \begin{tabular}{lccp{5.5cm}}
    \hline
    Variable & Mín & Máx & Observaciones \\
    \hline
    Resistencia observada (MPa) & 22{,}57 & 58{,}96 & P$_{25}$=31{,}19 · P$_{50}$=37{,}97 · P$_{75}$=43{,}34 MPa \\
    Temperatura (°C) & 10 & 32 & 3 niveles: 10\,°C (30{,}8\%), 25\,°C (46{,}2\%), 32\,°C (23{,}1\%) \\
    Humedad relativa (\%) & 20 & 70 & 3 niveles: 20\% (23{,}1\%), 50\% (46{,}2\%), 70\% (30{,}8\%) \\
    Relación a/c & 0{,}40 & 0{,}50 & 3 niveles: 0{,}40 (30{,}8\%), 0{,}45 (30{,}8\%), 0{,}50 (38{,}5\%) \\
    \hline
  \end{tabular}
\end{table}

\subsection{Metodología de validación}
Para evaluar la validez de las recomendaciones, se propone una validación cruzada o partición entrenamiento/validación con un conjunto de retención para verificación externa. Las métricas consideradas son:
\begin{itemize}
  \item Clasificación de aditivos (P0–PP3): exactitud, precisión, exhaustividad y F1 macro; matriz de confusión por clase.
  \item Proyección de resistencia (7/14/28 días): MAE, RMSE y MAPE por horizonte.
  \item Coherencia normativa: cumplimiento de restricciones y rangos indicativos según ACI 211.1.
  \item Trazabilidad: \% de recomendaciones con insumos/resultados/versionado correctamente almacenados en SQLite.
\end{itemize}

\begin{table}[H]
  \centering
  \caption{Métricas de clasificación de aditivos — vecino más cercano LOO (1\,872 registros, 7 clases)}
  \begin{tabular}{lcccc}
    \hline
    Clase & Precisión & Exhaustividad & F1 & Soporte \\
    \hline
    P0  & 0{,}537 & 0{,}528 & 0{,}532 & 468 \\
    P1  & 0{,}187 & 0{,}192 & 0{,}189 & 234 \\
    P2  & 0{,}156 & 0{,}132 & 0{,}143 & 234 \\
    P3  & 0{,}154 & 0{,}137 & 0{,}145 & 234 \\
    PP1 & 0{,}128 & 0{,}124 & 0{,}126 & 234 \\
    PP2 & 0{,}164 & 0{,}158 & 0{,}161 & 234 \\
    PP3 & 0{,}383 & 0{,}513 & 0{,}439 & 234 \\
    Macro-promedio & 0{,}244 & 0{,}255 & 0{,}248 & 1\,872 \\
    \hline
  \end{tabular}
\end{table}

\begin{table}[H]
  \centering
  \caption{Error de proyección de resistencia — validación LOO por horizonte}
  \begin{tabular}{lccc}
    \hline
    Horizonte & MAE (MPa) & RMSE (MPa) & MAPE (\%) \\
    \hline
    7 días  & 1{,}284 & 1{,}625 & 4{,}40 \\
    14 días & 1{,}767 & 2{,}207 & 4{,}61 \\
    28 días & 1{,}965 & 2{,}455 & 4{,}37 \\
    \hline
  \end{tabular}
\end{table}

\subsection{Resultados y análisis}
Se espera que DIAPCE proporcione: (i) sugerencias de dosificación base acordes con condiciones de entrada; (ii) clasificación P0–PP3 consistente con patrones observados; y (iii) proyecciones de resistencia razonables a 7/14/28 días. La interpretación debe considerar la cobertura de rangos: el desempeño es más confiable en regiones densamente muestreadas del conjunto de datos. Se recomienda reportar errores estratificados por tramos de temperatura, humedad y nivel de resistencia objetivo.

\subsection{Comparación con métodos tradicionales}
Para efectos académicos, se propone comparar las recomendaciones de DIAPCE con diseños elaborados mediante procedimiento manual guiado por ACI 211.1. Los indicadores comparativos incluyen: tiempo de decisión (flujo asistido vs. manual), coherencia con límites normativos, y diferencias en dosificación base y selección de aditivo. Cuando existan datos experimentales de laboratorio, debe contrastarse la proyección de resistencia con la resistencia ensayada.

\begin{figure}[H]
  \centering
  % \includegraphics[width=0.85\linewidth]{ruta/a/tu/figura.pdf}
  \fbox{\rule{0pt}{2in}\rule{0.9\linewidth}{0pt}}\\
  \caption{Comparación de resultados: DIAPCE vs. diseño manual ACI (placeholder)}
\end{figure}

\subsection{Sostenibilidad e impacto}
Aunque DIAPCE no implementa optimización automática de proporciones, su uso sistemático puede incentivar decisiones de dosificación más consistentes y selección adecuada de aditivos. Cuando se disponga de mediciones, puede evaluarse el impacto potencial en: (i) contenido de cemento frente a una línea base docente, y (ii) variación de resistencia a 28 días. Estos análisis deben interpretarse dentro de los márgenes normativos y del contexto de formación.

\subsection{Amenazas a la validez y limitaciones}
\begin{itemize}
  \item Sesgo de datos: el conjunto es regional; la generalización fuera del Meta debe justificarse.
  \item Medición: variabilidad experimental en laboratorio puede introducir ruido en las proyecciones a 7/14/28 días.
  \item Cobertura: desempeño reducido en combinaciones poco representadas (p. ej., extremos de temperatura/humedad).
  \item Alcance del sistema: la herramienta asiste decisiones; no sustituye criterio profesional ni verificación normativa.
\end{itemize}

\subsection{Reproducibilidad}
DIAPCE almacena localmente entradas y resultados (SQLite), lo que permite rastrear cada recomendación respecto a sus insumos y la versión del modelo. Se sugiere anexar en el repositorio tablas exportables (CSV) con métricas por cohorte de ensayos y capturas del flujo de uso para respaldar la discusión.
